\section{11 OpenCV}

\subsection{Basics \& Image I/O}

\mult{2}

\begin{concept}{OpenCV}
    Open-source computer-vision library (image processing, machine vision, ML).
    \begin{itemize}
        \item Intel 1999 (Bradsky); C++/Python/Java; CUDA/OpenCL GPU
        \item \important{OpenCV-Python} = the Python API
        \item images are \important{NumPy ndarrays}: 2-D (grayscale) or 3-D
              (\important{BGR} colour), \texttt{uint8} 0--255
    \end{itemize}
\end{concept}

\begin{definition}{HighGui (no display on target!)}
    HighGui (\texttt{imshow}, \texttt{namedWindow}, \texttt{createTrackbar},
    \texttt{waitKey}\dots) needs \important{X-Windows}, which is \important{not} in the
    MC2 Yocto image $\to$ use \important{\texttt{imwrite()} instead of \texttt{imshow()}}.
\end{definition}

\multend

\begin{KR}{Read / Write Images \& Pass Arguments}
\begin{lstlisting}[language=Python, style=basesmol, aboveskip=2pt, belowskip=2pt]
#!/usr/bin/env python3            # shebang: which interpreter
import cv2, sys, numpy as np
path = sys.argv[1]               # 1st command-line argument
img = cv2.imread(path, cv2.IMREAD_COLOR)   # ndarray, or None on failure
if img is None: sys.exit("Could not read the image.")
cv2.imwrite("out.png", img)      # type chosen from extension
\end{lstlisting}
    \texttt{imread} flag: \important{1 \texttt{IMREAD\_COLOR}} (default),
    \important{0 \texttt{IMREAD\_GRAYSCALE}}, \important{-1 \texttt{IMREAD\_UNCHANGED}}
    (incl. alpha). Run: \texttt{chmod +x f.py ; ./f.py img.jpg}.
\end{KR}

\begin{KR}{Capture Video}
\begin{lstlisting}[language=Python, style=basesmol, aboveskip=2pt, belowskip=2pt]
cap = cv2.VideoCapture(0)        # 0 = camera, or "file.webm"
if not cap.isOpened(): exit()
print(cap.get(cv2.CAP_PROP_FRAME_WIDTH))   # also FRAME_HEIGHT, FPS
while True:
ref, frame = cap.read()          # ref = success bool
gray = cv2.cvtColor(frame, cv2.COLOR_BGR2GRAY)
cv2.imwrite("gray.jpg", gray)
cap.release()
\end{lstlisting}
\end{KR}

\subsection{Core Operations}

\begin{KR}{Access \& Manipulate Pixels}
\begin{lstlisting}[language=Python, style=basesmol, aboveskip=2pt, belowskip=2pt]
(h, w, c) = img.shape[:3]        # rows, cols, channels
img.dtype                        # e.g. uint8
img.size                         # h*w*c
px = img[100, 100]               # [B, G, R] at y=100, x=100
blue = img[100, 100, 0]          # channel 0=B, 1=G, 2=R
img[100, 100, 2] = 0             # set red to 0
b, g, r = cv2.split(img)         # or faster: b = img[:,:,0]
img = cv2.merge((b, g, r))
img[:, :, 1] = 0                 # zero the green channel
\end{lstlisting}
    \important{Note:} indexing is \texttt{img[y, x]} (row, col); NumPy slicing is
    \important{faster} than \texttt{split}/\texttt{merge}.
\end{KR}

\mult{2}

\begin{definition}{Drawing Primitives}
    Each takes the image, geometry, \important{(B,G,R)} colour, then thickness
    (\important{-1 = fill}):
    \begin{itemize}
        \item \texttt{line(img, p0, p1, color, w)}
        \item \texttt{rectangle(img, c1, c2, color, t)}
        \item \texttt{circle(img, center, r, color, t)}
        \item \texttt{ellipse(img, c, (maj,min), angle, start, end, color, t)}
    \end{itemize}
\end{definition}

\begin{definition}{Borders \& Timing}
    \texttt{copyMakeBorder(src, top, bot, left, right, type, value)};
    type = \texttt{BORDER\_CONSTANT} / \texttt{REPLICATE} / \texttt{REFLECT}.
    Time a block:
\begin{lstlisting}[language=Python, style=basesmol, aboveskip=2pt, belowskip=1pt]
t1 = cv2.getTickCount()
# ... code ...
dt = (cv2.getTickCount()-t1) \
     / cv2.getTickFrequency()
\end{lstlisting}
\end{definition}

\multend

\begin{KR}{Region of Interest (ROI)}
    Slice \texttt{[start\_row:end\_row, start\_col:end\_col]} = \texttt{[y0:y1, x0:x1]}:
\begin{lstlisting}[language=Python, style=basesmol, aboveskip=2pt, belowskip=2pt]
roi = img[70:100, 150:190]       # cut out (30 x 40 region)
img[10:40, 10:50] = roi          # paste it elsewhere
\end{lstlisting}
\end{KR}

\subsection{Convolution}

%% >>> IMAGE (slide 32 / 35): kernel sliding over the image; border padding (3x3 kernel -> 1 line of padding) <<<

\begin{concept}{Convolution}
    \important{Element-wise multiply two matrices, then sum.} A small \important{kernel}
    slides left-to-right / top-to-bottom over the image; each output pixel = sum of
    (neighbourhood $\times$ kernel). Used for \important{blurring, sharpening, edge
    detection, Haar detection}.
\end{concept}

\mult{2}

\begin{formula}{Blur \& Sharpen Kernels}
    \textbf{Blur} (normalised average):
    $$K = \tfrac{1}{9}\begin{bmatrix}1&1&1\\1&1&1\\1&1&1\end{bmatrix}$$
    \textbf{Sharpen}:
    $$K = \begin{bmatrix}0&-1&0\\-1&5&-1\\0&-1&0\end{bmatrix}$$
\end{formula}

\begin{definition}{Padding \& Kernel Size}
    \begin{itemize}
        \item \important{Odd} kernel sizes preferred (have a center pixel)
        \item a 3$\times$3 kernel needs \important{1 line of padding}; generally
              \texttt{pad = (k-1)//2}
        \item padding methods: \important{replicate} (usual), zero, or wrap-around
    \end{itemize}
\end{definition}

\multend

\begin{KR}{Implement a Convolution}
\begin{lstlisting}[language=Python, style=basesmol, aboveskip=2pt, belowskip=2pt]
def convolve(image, kernel):
(iH, iW) = image.shape[:2]
(kH, kW) = kernel.shape[:2]
pad = (kW - 1) // 2                       # padding width
image = cv2.copyMakeBorder(image, pad, pad, pad, pad,
cv2.BORDER_REPLICATE)
output = np.zeros((iH, iW), dtype="float32")
for y in np.arange(pad, iH + pad):        # slide kernel
for x in np.arange(pad, iW + pad):
roi = image[y-pad:y+pad+1, x-pad:x+pad+1]  # neighbourhood
k = (roi * kernel).sum()              # multiply + sum
output[y-pad, x-pad] = k
return output
# build a 5x5 blur kernel and apply to a grayscale image:
kernel = np.ones((5, 5), dtype="float") * (1.0/(5*5))
gray = cv2.cvtColor(cv2.imread(path), cv2.COLOR_BGR2GRAY)
cv2.imwrite("out.jpg", convolve(gray, kernel))
\end{lstlisting}
\end{KR}

\begin{KR}{Built-in Filters (faster)}
\begin{lstlisting}[language=Python, style=basesmol, aboveskip=2pt, belowskip=2pt]
out = cv2.filter2D(gray, -1, kernel)   # ddepth -1 = same as src
g = cv2.GaussianBlur(img, (5,5), 0)    # ksize odd; sigma from size
a = cv2.blur(img, (5,5))               # average blur
m = cv2.medianBlur(img, 5)             # median blur
gray = cv2.cvtColor(img, cv2.COLOR_BGR2GRAY)   # also BGR2RGB, BGR2YUV
\end{lstlisting}
\end{KR}

\subsection{Object Detection (Haar Cascades)}

%% >>> IMAGE (slide 47): Haar-like features (edge, line, four-rectangle) <<<

\begin{concept}{Haar Cascade Classifiers}
    Object-detection using a \important{pretrained classifier in an XML file}.
    \begin{itemize}
        \item evaluate \important{Haar-like features} (edge/line/4-rectangle) on regions
        \item \important{AdaBoost} picks the most useful features, combining many
              \emph{weak} classifiers into one \emph{strong} classifier
        \item \important{cascade of stages} rejects obvious non-objects \important{early}
              $\to$ minimise the false-negative rate
    \end{itemize}
\end{concept}

\begin{KR}{Detect Faces / Objects}
\begin{lstlisting}[language=Python, style=basesmol, aboveskip=2pt, belowskip=2pt]
cascade = cv2.CascadeClassifier("haarcascade_frontalface_default.xml")
gray = cv2.cvtColor(image, cv2.COLOR_BGR2GRAY)   # must be grayscale
faces = cascade.detectMultiScale(gray, scaleFactor=1.1,
minNeighbors=3, minSize=(30,30), maxSize=(100,100),
flags=cv2.CASCADE_SCALE_IMAGE)
for (x, y, w, h) in faces:                       # returns rectangles
cv2.rectangle(image, (x,y), (x+w, y+h), (0,255,0), 2)
\end{lstlisting}
\end{KR}

\begin{definition}{detectMultiScale Parameters}
    \begin{itemize}
        \item \important{scaleFactor}: how much the image shrinks each pass.
              Small (e.g. 1.02) $\to$ fine search, accurate but \important{slow};
              large (1.5) $\to$ fast but misses
        \item \important{minNeighbors}: overlapping rectangles needed for a hit.
              Small $\to$ more \important{false positives}; large $\to$ more \important{misses}
        \item \important{minSize/maxSize}: window size range to detect
    \end{itemize}
\end{definition}

\subsection{OCR \& Codes}

\mult{2}

\begin{definition}{pytesseract (OCR)}
    Tesseract OCR engine: image $\to$ text (segmentation + recognition).
    \texttt{image\_to\_string(img, config=...)}.
    \begin{itemize}
        \item \important{PSM} (page seg): \texttt{6} single uniform block,
              \texttt{7} single line, \texttt{0} OSD only, \texttt{1} auto+OSD
        \item \important{OEM} (engine): \texttt{0} legacy, \texttt{1} LSTM neural net,
              \texttt{2} both, \texttt{3} default
    \end{itemize}
\end{definition}

\begin{KR}{Read QR / Barcodes (pyzbar)}
\begin{lstlisting}[language=Python, style=basesmol, aboveskip=2pt, belowskip=1pt]
from pyzbar import pyzbar
barcodes = pyzbar.decode(input)   # 1D + QR
for bc in barcodes:
x, y, w, h = bc.rect
info = bc.data.decode('utf-8')
cv2.rectangle(input, (x,y), (x+w,y+h), (0,255,0), 2)
\end{lstlisting}
\end{KR}

\multend

\begin{examplecode}{License-Plate OCR (slides 59--60)}
\begin{lstlisting}[language=Python, style=basesmol, aboveskip=2pt, belowskip=2pt]
import cv2, pytesseract, numpy as np
gray = cv2.cvtColor(cv2.imread("car.jpg"), cv2.COLOR_BGR2GRAY)
plate_cascade = cv2.CascadeClassifier(
"haarcascade_russian_plate_number.xml")
plates = plate_cascade.detectMultiScale(gray, scaleFactor=1.1,
minNeighbors=5, minSize=(30,30), flags=cv2.CASCADE_SCALE_IMAGE)
for (x, y, w, h) in plates:
roi = gray[y:y+h, x:x+w]
kernel = np.ones((5,5), np.uint8)
roi = cv2.erode(roi, kernel, iterations=1)     # clean up
text = pytesseract.image_to_string(roi, config=r'--oem 3 --psm 6')
\end{lstlisting}
\end{examplecode}

\raggedcolumns
